% Modernised reading — fonds Alexandre Grothendieck, shelfmark 153, pages 1-10.
%
% This is an INTERPRETATION, not a transcription. Everything the critical
% apparatus carried moves into footnotes. In any dispute of substance the
% transcription governs: whoever cites Grothendieck must cite batch-01.fr.tex.
%
% Pass: Opus 5 (claude-opus-5), 2026-09-13 — first pass, unchecked against the
% pages by a human. The folder was read whole; it holds one batch, pages 1-10.
% Pages 4, 6, 8, 10 are listings with nothing of his.
%
% What the folder is. The definition of an « analyseur » — a commutative
% pseudo-ring with a second, associative and unital, composition law, right
% composition being a ring homomorphism, and left composition with a sum or a
% product being expressible through finitely many elements — reached through
% two trial constructions (the axioms around a k_0-algebra, page 1; the
% binomial polynomials, page 2; the integer-valued polynomials of a pair
% k_0 ⊂ K_0, page 3), then the definition with its constants (page 5), and
% examples (pages 7, 9): all maps of a ring to itself, the polynomials, power
% series without constant term, the binomial analyseur, divided powers,
% lambda-rings. This is what is now called a plethory, or Tall-Wraith monoid.
%
% Notation fixed for the whole document. Omega an analyseur; o the
% composition, T its unit; Omega_0 the constants (page 5) and Omega^0 the
% F with F o 0 = 0. HOMONYM: page 1 writes Omega_0 for the object being
% axiomatised, which page 5 calls Omega, keeping Omega_0 for its constants;
% the reading writes Omega on page 1 and says so. F_n = binom(T, n). For a
% pair k_0 ⊂ K_0, Int(k_0) = {F in K_0[T] | F(k_0) ⊂ k_0} (his Omega).
%
% Substantive departures from the pages, each footnoted where it occurs:
% — page 3: K_0[X,Y] = K_0[X] (x)_{k_0} K_0[Y] holds when K_0 (x)_{k_0} K_0
%   = K_0 (e.g. K_0 a localisation of k_0); the uniqueness of the elements of
%   Omega (x) Omega representing F(x+y), F(xy) needs that and Omega free, the
%   latter being his own margin. The reading states both hypotheses.
% — page 5, 3) b): the sum is indexed by i where the family is indexed by
%   alpha.
% — page 7: Appl(A,A) satisfies 1) and 2) and not 3) in general (the leaf
%   says only « 1°) et 2°) »), nor does k[[T]]^+ satisfy 3); the reading
%   says so and gives the reason. The boxed « (Omega_A)_0 ≃ Gamma A ? » is
%   true, and the reading says why.
% — page 7: Omega -> Omega_{Omega_0} is not injective in general (k finite,
%   Omega = k[T]); the leaf's « on est content » is kept as a wish.
%
% Modern names given and footnoted as ours: plethory (Borger-Wieland, 2005),
% Tall-Wraith monoid (Tall-Wraith, 1970), biring and co-operations, ring of
% integer-valued polynomials Int(Z) and Pólya's basis, binomial ring,
% lambda-ring and plethysm on symmetric functions, divided powers.
%
% What the folder announces and never establishes: the axioms e) and f) of
% page 1, left without right-hand side; the analyseurs of divided powers and
% of lambda-rings, named on page 9 and not constructed; the meaning of W_k(A)
% on page 7, used and not defined.
\documentclass[11pt,a4paper]{article}
% Le préambule commun à toutes les transcriptions du fonds Grothendieck.
%
% Il tient en une idée : **l'incertitude se compose, elle ne se résout pas.**
% Une transcription qui lisse un mot douteux a détruit la seule chose pour
% laquelle on la faisait — un lecteur doit pouvoir savoir, à chaque endroit,
% ce qui était sur le feuillet et ce qui a été ajouté. Les six macros qui
% suivent sont donc l'appareil critique, et non de la décoration.
%
% Les mêmes commandes sont reconnues par `scripts/render.mjs`, qui produit la
% vue de lecture du site. Une macro ajoutée ici doit l'être aussi là-bas,
% faute de quoi le rendu échouera bruyamment — ce qui est voulu : un rendu
% silencieusement faux est pire qu'un rendu absent.

\ProvidesPackage{grothendieck}[2026/08/08 Transcriptions du fonds Grothendieck]

\usepackage{amsmath,amssymb,amsthm}
\usepackage{mathrsfs}
\usepackage{stmaryrd}
% Les diagrammes commutatifs sont partout dans ce fonds : c'est la langue
% même de la géométrie algébrique de Grothendieck, et une transcription qui
% les rendrait en prose ne transcrirait rien.
\usepackage{tikz-cd}
% Français par défaut — c'est la langue des feuillets — mais l'anglais est
% chargé aussi : la traduction et le résumé sont des documents anglais, et la
% typographie française (espaces avant les deux-points) y serait une faute.
% Ces documents basculent par \selectlanguage{english} en tête de corps.
\usepackage[english,french]{babel}
\usepackage{fontspec}
\usepackage{geometry}
\geometry{a4paper,margin=25mm,marginparwidth=28mm}
\usepackage{marginnote}
\usepackage{xcolor}
% ulem, not soul. `soul`'s \st and \ul are built on a character-by-character
% scan that XeLaTeX's font machinery breaks: the document fails with an
% undefined control sequence rather than degrading. ulem does the same two jobs
% and is Unicode-engine safe. `normalem` stops it hijacking \emph, which the
% transcriptions use for Grothendieck's own emphasis.
\usepackage[normalem]{ulem}
\usepackage{hyperref}
\hypersetup{colorlinks=true,linkcolor=black,urlcolor=black,pdfborder={0 0 0}}

% --- Métadonnées du lot -----------------------------------------------------
% Reprises telles quelles par le script de rendu, qui les affiche en tête de
% la vue de lecture. Elles fixent ce que le fichier prétend couvrir : sans
% elles, une transcription flotte sans point d'attache dans un fonds de seize
% mille feuillets.

\newcommand{\folder}[1]{\gdef\grothFolder{#1}}
\newcommand{\batch}[1]{\gdef\grothBatch{#1}}
\newcommand{\leaves}[2]{\gdef\grothFirst{#1}\gdef\grothLast{#2}}
\newcommand{\foldertitle}[1]{\gdef\grothFoldertitle{#1}}
\gdef\grothFolder{?}\gdef\grothBatch{?}\gdef\grothFirst{?}\gdef\grothLast{?}\gdef\grothFoldertitle{}

% --- Le résumé --------------------------------------------------------------
%
% Il ouvre la lecture modernisée, sous le titre « Résumé », et il est composé
% en petit. Ce n'est pas un ornement : c'est le seul passage du document qui
% ne soit pas des mathématiques, et un lecteur doit pouvoir le reconnaître
% avant de l'avoir lu — puis le sauter s'il connaît déjà le sujet, ou s'y
% arrêter s'il ne le connaît pas. Le filet qui le suit marque la reprise du
% corps.

\newenvironment{resume}
  {\small\setlength{\parskip}{0.45em}}
  {\par\medskip\hrule\medskip}

% --- Mots-clés ---------------------------------------------------------------
%
% En anglais, en clôture du résumé de la lecture modernisée : le vocabulaire
% moderne sous lequel chercher ce que les feuillets construisent. Le manifeste
% du site (`npm run manifest`) les extrait pour étiqueter la cote entière —
% cette ligne est la source unique des tags, il n'y a pas de fichier à côté.
\newcommand{\keywords}[1]{\par\smallskip\noindent
  {\footnotesize\textsc{Keywords} --- \textit{#1}}\par}

% --- L'appareil critique ----------------------------------------------------

% Le numéro de feuillet, en marge. C'est l'ancre qui permet de régler un
% désaccord sur une formule en regardant *un* feuillet plutôt qu'une chemise
% de six cents. Le site s'en sert aussi pour tourner les pages du fac-similé
% quand on fait défiler la transcription.
\newcommand{\leaf}[1]{%
  \marginnote{\footnotesize\color{gray!70}\textsf{#1}}%
  \label{leaf:#1}\ignorespaces}

% Illisible : on ne devine pas. Un mot inventé qui se lit comme les autres est
% le pire résultat possible.
\newcommand{\ill}{\textcolor{red!60!black}{[\,\ldots\,]}}

% Lecture douteuse : le mot est donné, et signalé comme incertain.
\newcommand{\uncertain}[1]{\uline{#1}}

% Ajout de l'éditeur — une lettre manquante, un mot sous-entendu.
\newcommand{\add}[1]{\textcolor{blue!50!black}{[#1]}}

% Biffé par Grothendieck lui-même. Ce qu'il a rayé fait partie du document :
% une rature dit souvent où la pensée a changé de direction.
\newcommand{\struck}[1]{\sout{#1}}

% Note du transcripteur, sur l'état matériel du feuillet.
\newcommand{\note}[1]{\par\smallskip{\small\color{gray!60!black}\textit{[#1]}}\par\smallskip}

% Note marginale de Grothendieck, distincte du corps du texte.
\newcommand{\marginal}[1]{\par\smallskip\noindent\hspace{1em}%
  {\small\color{gray!60!black}#1}\par\smallskip}

% --- Titre ------------------------------------------------------------------

\AtBeginDocument{%
  \begin{center}
    {\large\bfseries Fonds Alexandre Grothendieck --- cote n\textsuperscript{o}~\grothFolder}\\[2pt]
    {\itshape\grothFoldertitle}\\[4pt]
    {\small feuillets \grothFirst--\grothLast\ (lot \grothBatch)}\\[2pt]
    {\footnotesize\color{gray!60!black}Transcription établie d'après le fac-similé numérisé
     par l'Université de Montpellier.\\
     Les lectures incertaines sont soulignées, les illisibles notées [\,\ldots\,],
     les ajouts entre crochets.}
  \end{center}
  \vspace{1em}}

\endinput


\folder{153}
\pages{1}{10}
\dating{[à partir de 1982]}
\watermark{Édition de démonstration\\interprétation personnelle de l'œuvre}
\foldertitle{Analyseurs : notes manuscrites (s.d.) — lecture modernisée du dossier entier}

\begin{document}

\section*{Résumé}

\begin{resume}
Sur les nombres entiers, on sait additionner et multiplier. On sait aussi
appliquer certaines \emph{opérations} : élever au carré, prendre le coefficient
binomial $\binom{x}{n}$, qui est toujours un entier même si sa formule contient
une division. Ces opérations se composent entre elles, s'additionnent et se
multiplient, et elles forment à leur tour un anneau --- avec une loi de plus, la
composition.

Ces feuillets cherchent la bonne structure abstraite pour ce genre d'objet, et
lui donnent un nom : un \emph{analyseur}. C'est un anneau commutatif, pas
nécessairement unitaire, muni d'une seconde loi, la composition, associative
et unitaire, avec deux règles de compatibilité. La première dit que composer à
droite respecte l'addition et la multiplication : $(F + F') \circ G = F \circ G
+ F' \circ G$. La seconde est plus subtile : composer à gauche avec une somme ne
donne pas une somme --- $(x+y)^{2}$ n'est pas $x^{2} + y^{2}$ ---, mais il faut
qu'on puisse l'exprimer par un nombre fini d'opérations appliquées séparément
à $x$ et à $y$, comme $(x+y)^{2} = x^{2} + 2xy + y^{2}$.

L'exemple qui guide tout est celui des polynômes à coefficients rationnels qui
envoient les entiers sur des entiers --- les polynômes binomiaux $\binom{T}{n}$
et leurs combinaisons entières. La formule de Vandermonde,
$\binom{x+y}{n} = \sum_{i+j=n} \binom{x}{i}\binom{y}{j}$, est exactement la
seconde règle. D'autres exemples sont nommés au bout du dossier : l'analyseur des
puissances divisées, et celui des $\lambda$-anneaux, qui porte les opérations
de la théorie des représentations.

La notion est, à peu de chose près, celle que D.~Tall et G.~Wraith avaient
introduite en 1970, et qu'on étudie aujourd'hui sous le nom de
\emph{pléthorie} ; les feuillets ne citent personne. Ils s'attardent sur les
« constantes » d'un analyseur --- les opérations qui ne dépendent pas de leur
argument --- et font opérer l'analyseur sur elles, comme un anneau d'opérations
sur l'anneau qu'il analyse.

\keywords{plethory, Tall-Wraith monoid, biring, integer-valued polynomial, binomial ring, lambda-ring, divided powers, composition ring}
\end{resume}

\section*{Le fil du dossier, et les conventions}

\begin{enumerate}
\item[1.] \emph{Un premier jeu d'axiomes} (page 1).
\item[2.] \emph{Les polynômes binomiaux} (page 2).
\item[3.] \emph{Polynômes à valeurs entières relatifs à $k_{0} \subset K_{0}$}
(page 3).
\item[4.] \emph{Définition, constantes} (page 5).
\item[5.] \emph{Exemples} (pages 7 et 9).
\end{enumerate}

Les pages 1, 3, 5, 7 et 9 portent sa pagination 1 à 5 ; la page 2 est le verso
de la page 1.

\emph{Conventions.} $\Omega$ est un analyseur, $F \circ G$ sa composition, $T$
l'unité de la composition. $\Omega_{0}$ est l'ensemble de ses constantes, défini
à la page 5, et $\Omega^{0}$ celui des $F$ tels que $F \circ 0 = 0$. Pour
$n \in \mathbf{N}$,
$F_{n}(T) = \binom{T}{n} = \frac{T(T-1)\cdots(T-n+1)}{n!}$. Pour une inclusion
d'anneaux commutatifs $k_{0} \subset K_{0}$, on note
$\mathrm{Int}(k_{0}) = \{F \in K_{0}[T] \mid F(k_{0}) \subset k_{0}\}$, que la
page 3 appelle $\Omega$.

\emph{Une homonymie.} La page 1 appelle $\Omega_{0}$ l'objet même qu'elle
axiomatise, et la page 5 l'appelle $\Omega$, réservant $\Omega_{0}$ aux
constantes. On écrit $\Omega$ dans les deux cas.

\emph{Ce que le dossier annonce et n'établit pas} : les axiomes e) et f) de la
page 1, laissés sans second membre ; les analyseurs des puissances divisées et
des $\lambda$-anneaux, nommés et non construits ; l'objet $W_{k}(A)$ de la page
7, utilisé et non défini.

\pagerange{1}{1}
\section*{Un premier jeu d'axiomes (page 1)}

Soit $k_{0}$ un anneau de base --- le cas important est $k_{0} = \mathbf{Z}$ ---
et $\Omega$ une $k_{0}$-algèbre, éventuellement sans unité, contenant un élément
$T$, d'où $\varphi : k_{0}[T] \to \Omega$, et munie d'une opération
$(F, G) \mapsto F \circ G$ telle que :\footnote{La page écrit $\Omega_{0}$, qui
remplace un $k_{0}\langle T \rangle$ biffé trois fois ; voir la convention.}
\begin{enumerate}
\item[a)] $(F \circ G) \circ H = F \circ (G \circ H)$ ;
\item[b)] $T \circ F = F \circ T = F$ ;
\item[c)] $\varphi(F \circ G) = \varphi(F) \circ \varphi(G)$, la composition de
gauche étant la composition habituelle des polynômes ;
\item[d)] $F \mapsto F \circ G$ est un morphisme de $k_{0}$-algèbres ;
\item[e)] $F \circ (G + G') = \ ?$ ;
\item[f)] $F \circ (GG') = \ ?$ ;
\item[g)] $F \circ \lambda 1 = F(\lambda) \cdot 1$, où $\Omega$ opère sur $k_{0}$
par $(F, \lambda) \mapsto F(\lambda)$.
\end{enumerate}
Les deux axiomes e) et f), accolés, restent sans second membre : c'est
précisément ce que la page 5 saura écrire.\footnote{Une note en marge propose
« plutôt $W(\Omega) \times W(\Omega) \to W(\Omega)$ ». Au-dessus de d), un mot
d'une plume légère est lu « bilatérales » avec doute.}

La moitié inférieure du feuillet est un champ d'essais : une $k_{0}$-algèbre
commutative $A$ « sans unité », les lois ponctuelles
$(\lambda F)(x) = \lambda F(x)$, $(F + G)(x) = F(x) + G(x)$, $(FG)(x) =
F(x)G(x)$, la composition $(F \circ G)(x) = F(G(x))$, un morphisme
$W(\Omega) \to \mathrm{End}\, W(A)$, et, pour la première fois, l'idée de la
co-addition : $F(x + y) = \Delta F(x, y)$ dans $\Omega \otimes \Omega$.

\pagerange{2}{2}
\section*{Les polynômes binomiaux (page 2)}

Les $F_{n}$ sont dans $\mathbf{Q}[T]$ et envoient $\mathbf{Z}$ dans $\mathbf{Z}$.
On a
\[
F_{m} F_{n} = \sum_{p \leq m+n} c^{p}_{m,n} F_{p},
\qquad
F_{m} \circ F_{n} = \sum_{p \leq mn} d^{p}_{m,n} F_{p},
\qquad c^{p}_{m,n},\ d^{p}_{m,n} \in \mathbf{Z},
\]
et $\bigoplus_{n} \mathbf{Z} F_{n} \subset \mathbf{Q}[T]$ est une
$\mathbf{Z}$-algèbre commutative unitaire, stable par composition, engendrée
par les $F_{n}$ soumis aux relations $F_{0} = 1$ et $(\ast)$ ; on pose alors
$F_{1} = T$. De plus\footnote{Les entiers $c$, $d$, $\delta$ existent parce qu'un polynôme à une ou
deux variables qui prend des valeurs entières aux points entiers est
combinaison entière des $F_{n}$, ou des $F_{i}(X)F_{j}(Y)$ ; c'est le théorème de
G.~Pólya (1919), et $\bigoplus_{n} \mathbf{Z} F_{n}$ est l'anneau
$\mathrm{Int}(\mathbf{Z})$ des polynômes à valeurs entières. La page écrit la
co-addition sous la forme générale $\sum \gamma^{m}_{ij} F_{i}(P)F_{j}(Q)$ et
laisse inachevée la formule de la co-multiplication ; que $\gamma^{m}_{ij}$ vaille
1 si $i + j = m$ et 0 sinon est la formule de Vandermonde, qui n'est pas sur la
page. Elle écrit $\delta : \Omega \to \Omega \otimes_{\mathbf{Z}} \Omega$.}
\[
F_{m}(P + Q) = \sum_{i+j=m} F_{i}(P)\, F_{j}(Q),
\qquad
F_{m}(PQ) = \sum_{i,j} \delta^{m}_{ij}\, F_{i}(P)\, F_{j}(Q),
\qquad \delta^{m}_{ij} \in \mathbf{Z}.
\]

\pagerange{3}{3}
\section*{Polynômes à valeurs dans $k_{0}$ (page 3)}

Soit $k_{0} \subset K_{0}$ une inclusion d'anneaux commutatifs, et
\[
\Omega = \mathrm{Int}(k_{0}) = \{F \in K_{0}[T] \mid F(\lambda) \in k_{0}\
\forall \lambda \in k_{0}\}.
\]
C'est une sous-$k_{0}$-algèbre de $K_{0}[T]$, stable par composition, et elle
opère sur $k_{0}$ de sorte que $(F + G)(x) = F(x) + G(x)$, $(\lambda F)(x) =
\lambda F(x)$, $(FG)(x) = F(x)G(x)$ et $(F \circ G)(x) = F(G(x))$.

\textbf{Lemme.} Si $M$ est un $k_{0}$-module libre, l'ensemble des
$F \in K_{0}[T] \otimes_{k_{0}} M$ tels que $F(x) \in M$ pour tout $x \in k_{0}$
s'identifie à $\Omega \otimes_{k_{0}} M$ --- coordonnée par coordonnée dans une
base de $M$.\footnote{Le mot souligné qui ouvre l'énoncé est lu « Lemme »
sans certitude.}

\textbf{Les co-opérations.} Supposons $K_{0} \otimes_{k_{0}} K_{0} = K_{0}$ --- par
exemple $K_{0}$ localisé de $k_{0}$, comme $\mathbf{Q}$ de $\mathbf{Z}$ --- et
$\Omega$ libre comme $k_{0}$-module. Alors
$K_{0}[X, Y] = K_{0}[X] \otimes_{k_{0}} K_{0}[Y]$\footnote{La page écrit ces deux égalités avec un « ? » au-dessus du signe égal
et l'identité $K_{0}[X,Y] = K_{0}[X] \otimes_{k_{0}} K_{0}[Y]$ sans hypothèse ;
sa marge ajoute « OK si $\Omega$ est un $k_{0}$-module \emph{libre} ». La
seconde hypothèse est la sienne, la première la nôtre : sans elle, le produit
tensoriel sur $k_{0}$ est plus gros que $K_{0}[X,Y]$. Pour $\mathbf{Z} \subset
\mathbf{Q}$ les deux sont vraies, $\mathrm{Int}(\mathbf{Z})$ ayant pour base les
$F_{n}$. Le signe en exposant du second opérateur, lu $\times$ ou $+$, est rendu
ici par l'indice $m$.}, et, pour $F \in \Omega$, les
polynômes
\[
\Delta_{a} F(X, Y) = F(X + Y), \qquad \Delta_{m} F(X, Y) = F(XY)
\]
prennent des valeurs dans $k_{0}$ sur $k_{0} \times k_{0}$ ; par le lemme appliqué
à $M = \Omega$, ce sont des éléments bien déterminés de $\Omega \otimes_{k_{0}}
\Omega$,
\[
F(X + Y) = \sum_{i} G_{i}(X) H_{i}(Y),
\qquad
F(XY) = \sum_{\alpha} J_{\alpha}(X) K_{\alpha}(Y).
\]
Il en résulte les deux règles
\[
F \circ (P + Q) = \sum_{i} (G_{i} \circ P)(H_{i} \circ Q),
\qquad
F \circ (PQ) = \sum_{\alpha} (J_{\alpha} \circ P)(K_{\alpha} \circ Q),
\]
et, pour $\lambda \in k_{0}$, $F \circ \lambda T = \sum_{\alpha} J_{\alpha}(\lambda)
K_{\alpha} = \sum_{\alpha} K_{\alpha}(\lambda) J_{\alpha}$.

\pagerange{5}{5}
\section*{Définition, et constantes (page 5)}

\textbf{Définition.} Un \emph{analyseur} est un pseudo-anneau $\Omega$ --- un
groupe abélien muni d'une multiplication associative, commutative, biadditive,
pas nécessairement unitaire --- muni d'une loi de monoïde\footnote{En b), la page indexe la somme par $i$ et la famille par $\alpha$. Dans
la langue d'aujourd'hui, un analyseur est, à peu de chose près, un
\emph{monoïde de Tall-Wraith} (D.~Tall et G.~Wraith, 1970) ou une
\emph{pléthorie} au sens de J.~Borger et B.~Wieland (2005) : un anneau
commutatif muni d'une composition qui en fait un monoïde, les co-opérations
$\Delta_{a}$ et $\Delta_{m}$ étant des éléments de $\Omega \otimes \Omega$. La
condition 3) de la page 5 demande seulement l'existence d'une famille finie
qui convienne, là où la définition moderne se donne les co-opérations comme
structure --- ce que la page 3 montre bien déterminé dans le cas relatif. Les
noms ne sont pas sur la page, et rien ne dit que les feuillets connaissent
l'article de Tall et Wraith.}
$(F, G) \mapsto F \circ G$, d'unité notée $T$, pas nécessairement commutative,
telle que :
\begin{enumerate}
\item[1)] $+$ est associative et commutative, d'élément neutre $0$ ; $xy$ est
associative, commutative, biadditive ; $\circ$ est associative et unitaire ;
\item[2)] pour tout $G$, $F \mapsto F \circ G$ est un homomorphisme de
pseudo-anneaux :
$(F + F') \circ G = F \circ G + F' \circ G$ et $(FF') \circ G = (F \circ G)(F'
\circ G)$ ;
\item[3a)] (fonctorialité) pour tout $F$, il existe une famille finie $(F'_{i}, F''_{i})$ telle
que $F \circ (G' + G'') = \sum_{i} (F'_{i} \circ G')(F''_{i} \circ G'')$ pour
tous $G'$, $G''$ ;
\item[3b)] pour tout $F$, il existe une famille finie $(P'_{\alpha},
Q''_{\alpha})$ telle que $F \circ (G'G'') = \sum_{\alpha} (P'_{\alpha} \circ
G')(Q''_{\alpha} \circ G'')$ pour tous $G'$, $G''$.
\end{enumerate}

\textbf{N.B.} Que dire de $F \circ 0$ ? Par 2), $0 \circ G = 0$, donc
$(F \circ 0) \circ G = F \circ (0 \circ G) = F \circ 0$ : $F \circ 0$ est une
constante, au sens suivant.

\textbf{Définition.} $\lambda \in \Omega$ est une \emph{constante} si
$\lambda \circ F = \lambda$ pour tout $F$. La notion ne dépend que de la loi
$\circ$.

\textbf{Proposition.} L'ensemble $\Omega_{0}$ des constantes est un
sous-pseudo-anneau de $\Omega$, par 2). Pour $F \in \Omega$ et $\lambda \in
\Omega_{0}$, $F \circ \lambda \in \Omega_{0}$, car
$(F \circ \lambda) \circ G = F \circ (\lambda \circ G) = F \circ \lambda$. Ainsi
$\Omega$ opère sur $\Omega_{0}$ par
\[
(F, \lambda) \longmapsto F(\lambda) \overset{\text{déf}}{=} F \circ \lambda .
\]
Posant $\Omega^{0} = \{F \in \Omega \mid F(0) = 0\}$, on a
\[
\Omega = \Omega_{0} \oplus \Omega^{0},
\qquad F = F(0) + \bigl(F - F(0)\bigr),
\]
somme directe de groupes abéliens, $\Omega^{0}$ étant un idéal.\footnote{La page
énonce la décomposition sans la justifier. On a $(F - F(0)) \circ 0 = F(0) -
F(0) \circ 0 = 0$, $F(0)$ étant une constante, et une constante $\lambda$ telle
que $\lambda \circ 0 = 0$ est nulle ; $\Omega^{0}$ est un idéal parce que
$(FG) \circ 0 = (F \circ 0)(G \circ 0)$. Ces vérifications sont les nôtres. Le
premier adjectif de la proposition est illisible ; « sous-pseudo-anneau » est
ce que l'argument donne.}

\textbf{Définition.} L'analyseur est \emph{unitaire} si la multiplication a une
unité $1$ et que celle-ci est une constante : $1 \circ F = 1$ pour tout $F$.

\pagerange{7}{7}
\section*{Exemples (page 7)}

\textbf{Exemple 1.} Pour un pseudo-anneau commutatif $A$, $\Omega_{A} =
\mathrm{Appl}(A, A)$, avec les lois ponctuelles et la composition des
applications, satisfait 1) et 2), et il est unitaire si $A$ l'est. Ses constantes
sont les applications constantes, et $(\Omega_{A})_{0} \simeq A$. Il ne satisfait
pas 3) en général : pour $A$ infini, une application $(x, y) \mapsto F(x + y)$
n'est pas une somme finie de produits $F'(x)F''(y)$.\footnote{La page dit
« Les conditions 1°) et 2°) … vérifiées », sans parler de 3), et la
remarque sur 3) est la nôtre.}

\emph{Variante.} Si $A$ est une $k$-algèbre, on peut prendre
$\Omega_{A/k} = \mathrm{End}(W_{k}(A))$, avec $\Omega_{0} \simeq A$ ; et plus
généralement, pour un anneau commutatif $A$ d'un topos,
$\Omega = \mathrm{End}_{\text{faisceaux d'ens.}}(A)$, muni des lois $xy$ et
$\lambda \circ y$, unitaire si $A$ l'est. Il se demande, encadré :
$(\Omega_{A})_{0} \simeq \Gamma A$ ? C'est vrai : une constante $\lambda$ vérifie
$\lambda = \lambda \circ c_{0}$, où $c_{0}$ est l'endomorphisme constant de valeur
$0$, donc se factorise par l'objet final, et les constantes sont les sections
globales de $A$.\footnote{L'argument est le nôtre. La page n'écrit pas ce qu'est
$W_{k}(A)$ ; on le comprend comme le foncteur $R \mapsto R \otimes_{k} A$ sur les
$k$-algèbres, et c'est ce qui donne l'exemple suivant pour $A = k$. Plusieurs mots
de la page sont illisibles.}

\textbf{L'analyseur des polynômes.} Pour $A = k$, on trouve $\Omega \simeq k[T]$,
avec l'addition, la multiplication et la composition des polynômes : les
endomorphismes de la droite affine. Les conditions 3) sont vérifiées --- $F(X +
Y)$ et $F(XY)$ sont dans $k[X, Y] = k[X] \otimes_{k} k[Y]$ --- et l'anneau des
constantes est $k$.

\emph{Séries formelles.} On peut aussi prendre $k[[T]]^{+}$, les séries sans
terme constant, avec le produit et la composition ; les constantes s'y réduisent
à zéro. On peut encore admettre les séries $a_{0} + a_{1}T + \cdots$ dont le terme
constant est nilpotent, pour lesquelles la composition converge. Ce ne sont pas
des analyseurs au sens de 3).\footnote{La page ne dit pas si 3) est vérifiée. Elle
ne l'est pas : pour $F = \sum_{n} T^{n}$, si $F(X + Y)$ était une somme finie
$\sum_{i \leq N} a_{i}(X) b_{i}(Y)$, la matrice des coefficients de
$X^{p}Y^{q}$, qui est la matrice de Pascal $\binom{p+q}{p}$, serait de rang au
plus $N$ ; or ses mineurs principaux valent tous 1.}

\textbf{L'analyseur des constantes.} Pour un analyseur quelconque $\Omega$,
\[
\Omega \longrightarrow \mathrm{Appl}(\Omega_{0}, \Omega_{0}) = \Omega_{\Omega_{0}},
\qquad F \longmapsto (\lambda \mapsto F \circ \lambda),
\]
est un homomorphisme d'analyseurs, qui induit l'identité sur les constantes. « On
est content » quand il est injectif.\footnote{Il ne l'est pas toujours : pour
$k$ un corps fini à $q$ éléments et $\Omega = k[T]$, $T^{q} - T$ est envoyé sur
l'application nulle. Pour $\mathrm{Int}(\mathbf{Z})$ il l'est, un polynôme étant
déterminé par ses valeurs sur une partie infinie. L'exemple est le nôtre ; la
phrase de la page est à demi illisible.}

\pagerange{9}{9}
\section*{Exemples relatifs, et deux analyseurs nommés (page 9)}

\textbf{Exemple 2 (relatif).} Pour un couple $k_{0} \subset K_{0}$ d'anneaux
commutatifs, $\mathrm{Int}(k_{0}) \subset K_{0}[T]$, avec les lois induites par
celles de $K_{0}[T]$, est un analyseur, sous les hypothèses de la page 3 pour ce
qui est de 3). Le cas le plus intéressant est celui de $\mathbf{Z} \subset
\mathbf{Q}$ :

\textbf{L'analyseur binomial} $\mathrm{Int}(\mathbf{Z}) = \bigoplus_{n}
\mathbf{Z} F_{n}$, unitaire.\footnote{Les anneaux sur lesquels il opère sont les
\emph{anneaux binomiaux}, ceux où les $\binom{x}{n}$ ont un sens ; le nom n'est
pas sur la page.}

\textbf{Exemple 4.} L'analyseur des puissances divisées, non unitaire.

\textbf{Exemple 3.} L'analyseur des $\lambda$-anneaux, unitaire.

Ces deux derniers sont nommés et ne sont pas
construits.\footnote{Il les numérote dans l'ordre 4 puis 3, le premier chiffre
écrit par-dessus un autre. L'analyseur des $\lambda$-anneaux est aujourd'hui
l'anneau $\Lambda$ des fonctions symétriques muni du pléthysme, l'exemple
fondateur de Tall et Wraith ; celui des puissances divisées est la pléthorie
qui représente les anneaux à puissances divisées. Ces identifications sont les
nôtres.}

\end{document}
